% Tetrahedral numbers plus one that are squares
% Typesetting: LuaLaTeX (run lualatex two or three times).
\RequirePackage{fix-cm}
\documentclass[a4paper,11pt]{article}

\usepackage{fontspec}
\usepackage{newunicodechar}
\usepackage{amsmath,amssymb,amsthm,mathtools}
\usepackage{booktabs,array}
\usepackage{enumitem}
\usepackage[hidelinks]{hyperref}

% Unicode symbols in the Lean listing (Latin Modern Mono has no glyphs for them).
\newunicodechar{ℕ}{\ensuremath{\mathbb{N}}}
\newunicodechar{ℚ}{\ensuremath{\mathbb{Q}}}
\newunicodechar{∃}{\ensuremath{\exists}}
\newunicodechar{↔}{\ensuremath{\leftrightarrow}}
\newunicodechar{∨}{\ensuremath{\vee}}

\setlist{nosep,leftmargin=2em}
\setlist[enumerate]{label=(\arabic*)}
\newcolumntype{L}[1]{>{\raggedright\arraybackslash}p{#1}}

\theoremstyle{plain}
\newtheorem{theorem}{Theorem}[section]
\newtheorem{proposition}[theorem]{Proposition}
\newtheorem{lemma}[theorem]{Lemma}
\newtheorem{corollary}[theorem]{Corollary}
\theoremstyle{definition}
\newtheorem{definition}[theorem]{Definition}
\theoremstyle{remark}
\newtheorem{remark}[theorem]{Remark}
\numberwithin{equation}{section}
\numberwithin{table}{section}

\newcommand{\Z}{\mathbb{Z}}
\newcommand{\Q}{\mathbb{Q}}
\newcommand{\N}{\mathbb{N}}
\newcommand{\F}{\mathbb{F}}
\newcommand{\Rr}{\mathbb{R}}
\newcommand{\Nm}{\operatorname{N}}
\newcommand{\NF}{\operatorname{N}_F}
\newcommand{\tsec}[2]{\texorpdfstring{#1}{#2}}
\newcommand{\lean}[1]{\texttt{#1}}

\title{Tetrahedral numbers plus one that are squares}
\author{roiban1344}
\date{November 2026}

\begin{document}
\maketitle

\begin{abstract}
For a nonnegative integer $n$, the tetrahedral number $n(n+1)(n+2)/6$ plus $1$ is a square only for $n=0,5,8,22,24,63,13775$.
This equation defines an elliptic curve. Therefore, the integral points are usually determined by a computation based on lower bounds for linear forms in elliptic logarithms.
This paper gives a proof that does not use the theory of elliptic curves.
The central step of the proof combines the two equations into a single equation: the relative norm of an element $\xi$ of a quartic field is $11$.
We divide $\xi$ by $11$ and reduce it to a power of a unit. We determine the exponent by the $2$-adic uniqueness lemma.
The only numerical computations are products of particular elements and congruences modulo small numbers; no search for candidate solutions is made.
The whole proof is formalized in Lean~4 and Mathlib.
\end{abstract}

\tableofcontents

%======================================================================
\section{Introduction}
%======================================================================

\subsection{Main theorem}

Throughout, $\N=\{0,1,2,\dots\}$, and $T(n)=n(n+1)(n+2)/6$ is the tetrahedral number (OEIS A000292).

\begin{theorem}\label{thm:main}
For $n\in\N$, there is $m\in\N$ with $T(n)+1=m^2$ if and only if
\[
n\in\{0,\ 5,\ 8,\ 22,\ 24,\ 63,\ 13775\}.
\]
The corresponding values of $m$ are $1,6,11,45,51,209,660099$, in this order.
\end{theorem}

The sequence of the $m$ is registered in the OEIS as A216269, ``$k$ such that $k^2-1$ is a tetrahedral number'', with these seven terms.

\subsection{Relation to elliptic curves}

Multiplying the equation $T(n)+1=m^2$ by $6$ gives $6m^2=n(n+1)(n+2)+6$.
With $X=6n$ and $Y=36m$, this becomes
\[
Y^2=X^3+18X^2+72X+1296.
\]
This is a problem about the integral points of an elliptic curve.
Hence the finiteness of the set of solutions follows from Siegel's theorem.
However, Siegel's theorem gives no upper bound for the size of the solutions.
Therefore, to find all the integral points, one usually uses the elliptic logarithm method~\cite{ST94,GPZ94}.
In this method, one finds a basis of the Mordell--Weil group, obtains an upper bound for the coefficients from a lower bound for linear forms in elliptic logarithms, and finally examines the remaining range by lattice reduction.

The proof in this paper uses none of these.
It uses only the following.
\begin{itemize}
\item Computations in five rings, each a free $\Z$-module of rank $4$ with an explicit multiplication.
\item Inequalities between real numbers.
\item Congruences and divisibility of integers.
\end{itemize}
We do not use class numbers, the unit theorem, unique factorization of ideals, or $p$-adic analysis.
For the units, we show directly by real inequalities that one explicit element generates all relative units (Section~\ref{sec:units}).
To determine the exponent, we use an argument by repeated divisibility by powers of $2$ instead of the theory of $p$-adic zeros (Section~\ref{sec:dyadic}).

\subsection{Structure of the proof}

The proof proceeds in the following order.
\begin{enumerate}
\item Split the equation into five systems $n+3=dr^2$, $n^2+2=ey^2$ (Section~\ref{sec:branches}).
\item Write the solutions of the second equation in the form of a square in two integers (Section~\ref{sec:parameters}).
\item Combine the two equations into an element $\xi$ of relative norm $11$ (Section~\ref{sec:norm11}).
\item By division by $11$, write $\xi$ as the product of one of finitely many initial elements and a relative unit (Section~\ref{sec:division}).
\item Show that the relative units are exhausted by the powers of one element $\rho$ and signs (Section~\ref{sec:units}).
\item Remove the negative sign and the unneeded initial element (Section~\ref{sec:signs}).
\item Determine the exponent of $\rho$ by the $2$-adic uniqueness lemma (Sections~\ref{sec:dyadic} and~\ref{sec:exponents}).
\item Read off $n$ from the coordinates of the resulting element (Section~\ref{sec:conclusion}).
\end{enumerate}
In the last step, the $0$th coordinate of $\xi$ equals $n+3$ from the start. Therefore we can read off $n$ without passing through other quantities.

\subsection{Formalization}

Every step is formalized in Lean~4 and Mathlib.
The final theorem has the form shown in Appendix~\ref{app:lean}, and the Lean kernel checks it.
It depends only on the three standard axioms of Lean (\lean{propext}, \lean{Classical.choice}, \lean{Quot.sound}).
We do not use \lean{sorry} or \lean{native\_decide}.
The kernel checks the products of particular elements and the congruences modulo small numbers by \lean{decide}.
The correspondence between the steps of the text and the Lean files is summarized in the table in Appendix~\ref{app:lean}.

\subsection{Notation}

Let $e$ be a positive integer that is not a square, and write $\omega=\sqrt e$.
The norm on $\Z[\omega]=\{a+b\omega : a,b\in\Z\}$ is $\NF(a+b\omega)=a^2-eb^2$.
For an element $\alpha$ of $\Z[\omega]$, we write $\alpha'$ for its conjugate $a-b\omega$.

%======================================================================
\section{Five systems of equations}\label{sec:branches}
%======================================================================

Multiplying $T(n)+1=m^2$ by $6$ and rearranging, we get
\begin{equation}\label{eq:shift}
6m^2=(n+3)(n^2+2).
\end{equation}

\begin{proposition}\label{prop:branches}
If $(n,m)$ is a solution of \eqref{eq:shift} in natural numbers, then there are positive integers $r,y$ such that
\begin{equation}\label{eq:pair}
n+3=dr^2,\qquad n^2+2=ey^2
\end{equation}
holds for one of the five pairs $(d,e)$ in Table~\ref{tab:branches}.
\end{proposition}

\begin{table}[htbp]
\centering
\caption{The five cases}\label{tab:branches}
\begin{tabular}{c|ccccc}
\toprule
$d$ & $1$ & $2$ & $3$ & $11$ & $66$\\
$e$ & $6$ & $3$ & $2$ & $66$ & $11$\\
\bottomrule
\end{tabular}
\end{table}

\begin{proof}
Write $n+3=dr^2$ with $d$ a squarefree positive integer.
First, every common divisor of $n+3$ and $n^2+2$ divides $11$. In fact, $(n^2+2)-(n-3)(n+3)=11$.
Next, comparing the parity of the exponents of the prime factors on both sides of \eqref{eq:shift}, we see that $d$ divides $6(n^2+2)$.
Therefore, if a prime $p$ dividing $d$ does not divide $6$, then $p$ is a common divisor of $n+3$ and $n^2+2$, and so $p=11$.
Hence $d\mid66$.

Also, by \eqref{eq:shift}, $6d(n^2+2)$ is a square.
Write $6d=k^2e$ with $e$ squarefree. Then $e(n^2+2)$ is also a square, and we get $n^2+2=ey^2$.
For $d=1,2,3,6,11,22,33,66$ we have $e=6,3,2,1,66,33,22,11$.

It remains to exclude $d=6,22,33$.
For $d=6,22$ we have $e=1,33$. In both cases $e\equiv1\pmod4$, so $n^2+2\equiv y^2\pmod4$. Squares are $0$ or $1$ modulo $4$, so this is impossible.
For $d=33$, $n+3=33r^2$ gives $3\mid n$. Then $n^2+2=22y^2$ gives $22y^2\equiv2\pmod3$, that is, $y^2\equiv2\pmod3$. This is also impossible.
\end{proof}

Conversely, a solution of \eqref{eq:pair} gives a solution of \eqref{eq:shift}. In fact, $de=6k'^2$ ($k'=1,1,1,11,11$), so we may take $m=k'ry$.

%======================================================================
\section{Quadratic square representations}\label{sec:parameters}
%======================================================================

In this section we express the solutions of the second equation $n^2+2=ey^2$ by quadratic forms in two integers $a,b$.
We do not use the general theory of solutions of Pell equations. We construct the parameters by a descent: an argument that repeatedly produces a smaller solution from a given one until it reaches the smallest solution, in the manner of Fermat's infinite descent.

\subsection{Descent}

In the four cases with $e\ne2$, we can write $e=a_0^2+2$ ($a_0=2,1,8,3$).
Put $c_0=a_0^2+1$. Then $c_0+a_0\omega$ has norm $1$.

\begin{lemma}\label{lem:descent}
Suppose that $n\ge0$ and $y>1$ satisfy $n^2+2=ey^2$. Then
\[
N=c_0n-a_0ey,\qquad Z=c_0y-a_0n
\]
satisfy $0\le N<n$, $Z>0$ and $N^2+2=eZ^2$. Moreover, $n=c_0N+a_0eZ$ and $y=c_0Z+a_0N$.
\end{lemma}

\begin{proof}
We have $N+Z\omega=(n+y\omega)(c_0-a_0\omega)$, so the norm is preserved. The last two equations come from multiplying back by $c_0+a_0\omega$.
We show the inequalities in the following order. Note that $e=c_0+1$.
\begin{enumerate}
\item $a_0y<n$. Otherwise $n^2\le a_0^2y^2=ey^2-2y^2<ey^2-2=n^2$ (using $y>1$).
\item $Z>0$. This follows from the identity $(c_0y)^2-(a_0n)^2=y^2+2a_0^2>0$.
\item $|N|<n$. The inequality $N>-n$ is equivalent to $(c_0+1)n>a_0ey$, that is, to $n>a_0y$. The inequality $N<n$ is equivalent to $a_0n<ey$, and by (2) we have $a_0n<c_0y<ey$.
\item $a_0Z\le|N|$. Apply the same argument as in (1) to $(|N|,Z)$ (using $Z\ge1$).
\item $N\ge0$. Suppose $N<0$. Then $n=c_0N+a_0eZ=a_0eZ-c_0|N|\le a_0eZ-c_0a_0Z=a_0Z\le|N|<n$, a contradiction.
\end{enumerate}
\end{proof}

This descent stops at $y=1$, and then $n=a_0$.

\begin{proposition}\label{prop:parameters}
Let $e=A^2+2$, and put $Q_A(a,b)=a^2-2b^2+2Aab$.
If $n\ge0$ and $y>0$ satisfy $n^2+2=ey^2$, then there are $A=a_0$ or $A=-a_0$ and integers $a,b$ such that
\begin{equation}\label{eq:param}
Q_A(a,b)=1,\qquad n=A(a^2-2b^2)-4ab.
\end{equation}
\end{proposition}

\begin{proof}
We use induction on $n$. If $y=1$, then $n=a_0$, and $(A,a,b)=(a_0,1,0)$ works.
If $y>1$, the pair $(N,Z)$ of Lemma~\ref{lem:descent} satisfies $0\le N<n$, so it has a representation by the induction hypothesis. From it, the following update gives a representation of $(n,y)$:
\[
(a_0,a,b)\mapsto(-a_0,\ a,\ b-a_0a),\qquad(-a_0,a,b)\mapsto(a_0,\ a-2a_0b,\ b).
\]
One checks by expansion that \eqref{eq:param} is preserved.
\end{proof}

\subsection{Choice of sign}

Substituting $a^2-2b^2=1-2Aab$ from the first equation of \eqref{eq:param} into the second, we get
\begin{equation}\label{eq:offset}
n=A-2(A^2+2)ab.
\end{equation}
Combining this with the first equation $n+3=dr^2$, one of the two signs of $A$ is excluded in each case.
\begin{itemize}
\item $d=1$, $A=2$: reading $Q_2(a,b)=1$ modulo $4$, $b$ is even. Then $r^2=n+3=5-12ab\equiv5\pmod8$, which is impossible.
\item $d=2$, $A=1$: $2r^2=4-6ab$ gives $r^2\equiv2\pmod3$, which is impossible.
\item $d=11$, $A=-8$: $11r^2=-5-132ab$ gives $11\mid5$, which is impossible.
\item $d=66$, $A=3$: $66r^2=6-22ab$ gives $11\mid6$, which is impossible.
\end{itemize}
Therefore $A=-2,\,-1,\,8,\,-3$ (for $d=1,2,11,66$).

\subsection{The case \tsec{$e=2$}{e=2}}

For $d=3$, $e=2$, the equation $n+3=3r^2$ gives $3\mid n$. Moreover, $n$ is even (by $n^2+2=2y^2$).
Write $n=2N$. Then $y^2-2N^2=1$.
Modulo $4$, $y$ is odd and $N$ is even.
The product of the coprime consecutive integers $(y-1)/2$ and $(y+1)/2$ equals $N^2/2$. Hence each factor is a square or twice a square.
From this we get an element $a+b\sqrt2$ of norm $\pm1$ with $y+N\sqrt2=(a+b\sqrt2)^2$.
Then $n=2N=4ab$.
If the norm were $-1$, then $a^2+b^2\equiv2\pmod3$ would give $3\nmid ab$, contrary to $3\mid n$.
Therefore the norm is $1$, and changing the sign of $b$ gives the representation \eqref{eq:param} with $A=0$.

\subsection{Square representation}

\begin{proposition}\label{prop:square}
If $n^2+2=ey^2$, $y>0$, and \eqref{eq:param} holds, then $y=a^2+2b^2$ and
\[
n+y\omega=(A+\omega)\bigl(a+(A-\omega)b\bigr)^2.
\]
\end{proposition}

\begin{proof}
Substituting \eqref{eq:param} into the identity
\[
\bigl(A(a^2-2b^2)-4ab\bigr)^2+2Q_A(a,b)^2=(A^2+2)(a^2+2b^2)^2,
\]
we get $n^2+2=e(a^2+2b^2)^2$. Comparing with $y>0$, we get $y=a^2+2b^2$. The second part follows by expansion.
\end{proof}

To match the rings used later, we multiply once by a known unit and replace the fixed element.
The result has the form $w=n+y\omega=\beta\eta^2$ with $\eta=\eta_0+\eta_1\omega$. The values of $\beta,\eta_0,\eta_1$ in each case are given in Table~\ref{tab:beta}.
Each row of the table is a quadratic identity in $a,b$ and can be checked by expansion.

\begin{table}[htbp]
\centering
\caption{The square representation $n+y\omega=\beta(\eta_0+\eta_1\omega)^2$}\label{tab:beta}
\begin{tabular}{ccccc}
\toprule
$d$ & $A$ & $\beta$ & $\eta_0$ & $\eta_1$\\
\midrule
$1$ & $-2$ & $22+9\omega$ & $5a+2b$ & $-2a-b$\\
$2$ & $-1$ & $5+3\omega$ & $2a+b$ & $-a-b$\\
$3$ & $0$ & $\omega$ & $a$ & $-b$\\
$11$ & $8$ & $8+\omega$ & $a+8b$ & $-b$\\
$66$ & $-3$ & $63+19\omega$ & $10a+3b$ & $-3a-b$\\
\bottomrule
\end{tabular}
\end{table}

%======================================================================
\section{Elements of relative norm \tsec{$11$}{11}}\label{sec:norm11}
%======================================================================

\subsection{The idea}

Let $F=\Q(\omega)$ and consider a quadratic extension $L=F(\theta)$ with $\theta^2\in F$.
The two equations \eqref{eq:pair} become one equation: the relative norm $\Nm_{L/F}(\xi)$ of an element $\xi\in L$ is $11$ (Proposition~\ref{prop:xi}).
An element of a fixed relative norm is a product of one of a few elements and a unit of relative norm $1$.
Sections~\ref{sec:division} and~\ref{sec:units} carry out this decomposition explicitly, as a division by $11$ and a determination of the units.

\subsection{Five rings}

In each case we define a commutative ring $R=R_d$ by putting a multiplication on $\Z^4$.
The ring $R$ contains $\Z[\omega]$ as a subring and has an element $\theta$ with $\theta^2=\Delta\in\Z[\omega]$. It is a subring of $L=F(\theta)$.
For $d\ne11$, $R=\Z[\omega]\oplus\Z[\omega]\theta$. For $d=11$, $R=\Z[\omega]\oplus(\Z+\Z\,\omega/11)\theta$.
The value of $\Delta$ and the meaning of the coordinates are given in Table~\ref{tab:rings}.

\begin{table}[htbp]
\centering
\caption{The five rings. The element represented by $x=(x_0,x_1,x_2,x_3)$, and $\theta^2=\Delta$}\label{tab:rings}
\begin{tabular}{cccc}
\toprule
$d$ & $e$ & element represented by $x$ & $\Delta=\theta^2$\\
\midrule
$1$ & $6$ & $(x_0+x_1\omega)+(x_2+x_3\omega)\theta$ & $44+18\omega$\\
$2$ & $3$ & same as above & $20+12\omega$\\
$3$ & $2$ & same as above & $6\omega$\\
$11$ & $66$ & $(x_0+x_1\omega)+(x_2+x_3\omega/11)\theta$ & $176+22\omega$\\
$66$ & $11$ & $(x_0+x_1\omega)+(x_2+x_3\omega)\theta$ & $756+228\omega$\\
\bottomrule
\end{tabular}
\end{table}

For $d=1,2,3,11$ we have $\Delta=2d\beta$. For $d=66$ we have $2d\beta=11\Delta=(\omega\theta)^2$.
In what follows, put $s=\theta$ ($d\ne66$) and $s=\omega\theta$ ($d=66$). In every case $s^2=2d\beta$.

\begin{definition}
\begin{enumerate}
\item Let $\tau$, the \textbf{relative conjugation}, be the restriction to $R$ of the conjugation of $L/F$ ($\theta\mapsto-\theta$). In coordinates, $\tau(x)=(x_0,x_1,-x_2,-x_3)$.
\item We write $\Nm(x)=x\,\tau(x)$. This is the relative norm $\Nm_{L/F}(x)$; for $x\in R$ it lies in $\Z[\omega]$.
\item An element $x$ with $\Nm(x)=1$ is called a \textbf{relative unit}. A relative unit $x$ is invertible, and $x^{-1}=\tau(x)$.
\item We write the first two coordinates of $x$ as $P(x)=x_0$, $Q(x)=x_1$, and put $\pi(x)=x_0+x_1\omega\in\Z[\omega]$.
\end{enumerate}
\end{definition}

The multiplication table of each ring is determined by the relations in Table~\ref{tab:rings}.
In Lean, these are defined as multiplications on $\Z^4$, and the axioms of a commutative ring are checked.
(Whether $R$ is the full ring of integers does not matter for the argument.)

\subsection{Elements of relative norm \tsec{$11$}{11}}

\begin{proposition}\label{prop:xi}
For a solution of \eqref{eq:pair}, put $\xi=w+3+r\eta s\in R$, where $w=n+y\omega$. Then
\begin{equation}\label{eq:norm11}
\Nm(\xi)=11.
\end{equation}
Moreover,
\begin{equation}\label{eq:f}
f(\xi)=0,\qquad P(\xi)=n+3=dr^2,
\end{equation}
where $f(x)=(P(x)-3)^2-eQ(x)^2+2$.
\end{proposition}

\begin{proof}
Note that $\NF(w)=n^2-ey^2=-2$. Then
\begin{align*}
\xi\,\tau(\xi)&=(w+3)^2-r^2\eta^2s^2=(w+3)^2-2dr^2\beta\eta^2\\
&=(w+3)^2-2(n+3)w=w^2-2nw+9.
\end{align*}
Moreover, $w+w'=2n$ and $ww'=-2$, so $w^2-2nw=-ww'=2$. Hence $\Nm(\xi)=11$.
Since $\pi(\xi)=w+3$, we have $P(\xi)=n+3$ and $Q(\xi)=y$, and $f(\xi)=n^2-ey^2+2=0$.
\end{proof}

Here $11=3^2+2$. The two equations \eqref{eq:pair} combine into one small relative norm.
The coordinates of $\xi$ are $(n+3,y,r\eta_0,r\eta_1)$ for $d=1,2,3$, $(n+3,y,r\eta_0,11r\eta_1)$ for $d=11$, and $(n+3,y,11r\eta_1,r\eta_0)$ for $d=66$.

The function $f$ depends only on the first two coordinates. Therefore $f(\tau(x))=f(x)$ and $P(\tau(x))=P(x)$.

%======================================================================
\section{Division by \tsec{$11$}{11}}\label{sec:division}
%======================================================================

We call the elements in Table~\ref{tab:initial} the \textbf{initial elements}. Each of them has relative norm $11$.

\begin{table}[htbp]
\centering
\caption{The initial elements and the relative unit $\rho$}\label{tab:initial}
\begin{tabular}{cccc}
\toprule
$d$ & $\xi_0$ & $\xi_1$ & $\rho$\\
\midrule
$1$ & $(1,1,5,-2)$ & --- & $(9,4,2,0)$\\
$2$ & $(2,1,2,-1)$ & $(2,3,1,0)$ & $(3,2,1,0)$\\
$3$ & $(3,1,1,0)$ & --- & $(17,12,6,4)$\\
$11$ & $(11,1,1,0)$ & --- & $(33,4,2,2)$\\
$66$ & $(0,1,0,0)$ & --- & $(19,6,1,0)$\\
\bottomrule
\end{tabular}
\end{table}

\begin{proposition}\label{prop:division}
If $\Nm(\xi)=11$, then there are an initial element $p$ and a relative unit $u$ such that $\xi=u\,p$ or $\xi=u\,\tau(p)$.
\end{proposition}

\begin{proof}
In each case it suffices to show that all four coordinates of $\xi\,\tau(p)$ or of $\xi\,p$ are divisible by $11$.
In fact, if so, put $u=\xi\tau(p)/11$. Then $\Nm(u)=\Nm(\xi)\Nm(p)/121=1$, and $\xi\,\tau(p)\,p=11\xi$ gives $\xi=up$.

For $d=1,3$, $e=6,2$ is not a quadratic residue modulo $11$.
Therefore $\Z[\omega]/(11)$ is the field $\F_{121}$ with $121$ elements.
Write $\xi_0=\alpha_0+\beta_0\theta$. Then $\Nm(\xi_0)=\alpha_0^2-\Delta\beta_0^2=11$, and $\beta_0$ is coprime to $11$.
Therefore $\Delta$ has square roots $\pm\delta$ in $\F_{121}$ ($\delta=\alpha_0/\beta_0$).
Then $\theta\mapsto\pm\delta$ defines two ring homomorphisms $\chi_\pm:R\to\F_{121}$, and $\chi_-=\chi_+\circ\tau$.
Since $\chi_+(\xi)\chi_-(\xi)=\chi_+(\Nm(\xi))=0$, $\xi$ vanishes under one of them.
The elements $\xi_0$ and $\tau(\xi_0)$ vanish on different sides. Let $p$ be the one that vanishes on the same side as $\xi$. Then $\xi\tau(p)$ vanishes on both sides.
The intersection of the two kernels is $11R$.

For $d=2$, $e=3$ is a quadratic residue modulo $11$, and $\Z[\omega]/(11)\cong\F_{11}\times\F_{11}$.
In each of the two components there are two choices for the value of $\theta$ (the existence of square roots of $\Delta$ follows from $\xi_0$ as above).
We call the pair recording, in each component, the value at which $\xi$ vanishes the \textbf{type} of $\xi$. There are four types.
The elements $\xi_0,\xi_1,\tau(\xi_0),\tau(\xi_1)$ realize these four types, one each. Hence, by the same argument as above, $\xi\tau(p)\in11R$ for the $p$ with the same type as $\xi$.

For $d=11$, reading the relative norm modulo $11$ gives $x_0^2-96x_3^2\equiv0$.
Since $96\equiv8$ is not a quadratic residue modulo $11$, we get $11\mid x_0$ and $11\mid x_3$.
From this, $\xi\tau(\xi_0)\in11R$ follows.

For $d=66$, similarly, $x_0^2-756x_2^2\equiv0\pmod{11}$ and $756\equiv8$ give $11\mid x_0$ and $11\mid x_2$.
From this, $\xi\xi_0\in11R$ follows. Note that $\tau(\xi_0)=\xi_0$.
\end{proof}

%======================================================================
\section{Relative units}\label{sec:units}
%======================================================================

\begin{proposition}\label{prop:units}
In each case, every relative unit has the form $\pm\rho^m$ ($m\in\Z$), where $\rho$ is the element in Table~\ref{tab:initial}.
\end{proposition}

A direct computation shows that $\rho$ is a relative unit. What we must show is that there are no other relative units.

\subsection{Two real embeddings}

For $\Delta=\Delta_0+\Delta_1\omega$, we have $\sigma(\Delta)=\Delta_0+\Delta_1\sqrt e>0$ and $\sigma'(\Delta)=\Delta_0-\Delta_1\sqrt e<0$.
Put $D_+=\sigma(\Delta)$ and $D_-=-\sigma'(\Delta)$.
Define $\sigma:R\to\Rr$ by $\omega\mapsto\sqrt e$, $\theta\mapsto\sqrt{D_+}$.
We write $\alpha\mapsto\alpha'$ for the conjugation of $\Z[\omega]$ defined by $\omega\mapsto-\sqrt e$.

Write a relative unit as $v=\alpha+\gamma\theta$, where $\alpha\in\Z[\omega]$ and $\gamma\in\Z[\omega]$ ($\gamma\in\Z+\Z\,\omega/11$ for $d=11$).
Taking the conjugate of $\Nm(v)=\alpha^2-\Delta\gamma^2=1$, we get $\alpha'^2+D_-\gamma'^2=1$. Therefore $|\gamma'|\le D_-^{-1/2}$.

\begin{lemma}\label{lem:gap}
If $\gamma\ne0$, then
\begin{equation}\label{eq:gap}
\bigl(|\sigma(v)|-|\sigma(v)|^{-1}\bigr)^2\ge4M,\qquad M=\nu^2D_+D_-,
\end{equation}
where $\nu=1$ ($d\ne11$) and $\nu=3/11$ ($d=11$).
\end{lemma}

\begin{proof}
Since $\sigma(\tau(v))=\sigma(v)^{-1}$, we have $\sigma(v)-\sigma(v)^{-1}=2\sigma(\gamma)\sqrt{D_+}$.
If $\gamma\ne0$, then $|\sigma(\gamma)\gamma'|=|\NF(\gamma)|\ge\nu$.
In fact, for $d\ne11$, $\NF(\gamma)$ is a nonzero integer.
For $d=11$, write $\gamma=c+k\omega/11$. Then $\NF(\gamma)=(11c^2-6k^2)/11$, and the cases where the numerator has absolute value $1$ or $2$ are excluded modulo $3$ and modulo $11$.
Therefore
\begin{equation}\label{eq:gamma}
\sigma(\gamma)^2\ge\nu^2/\gamma'^2\ge\nu^2D_-,
\end{equation}
and we get \eqref{eq:gap}.
\end{proof}

The lemma says that the real value $|\sigma(v)|$ of a relative unit $v\ne\pm1$ stays a fixed distance away from $1$.
The values of $4M$ and of $\Lambda=\sigma(\rho)$ are given in Table~\ref{tab:gap}. For $\Lambda$, we use only the upper bound in the table and $\Lambda>1$.

\begin{table}[htbp]
\centering
\caption{The constant $4M$ of Lemma~\ref{lem:gap} and upper bounds for $\Lambda=\sigma(\rho)$}\label{tab:gap}
\begin{tabular}{cccccc}
\toprule
$d$ & $1$ & $2$ & $3$ & $11$ & $66$\\
\midrule
$4M$ & $32$ & $128$ & $288$ & $288$ & $1152$\\
$\Lambda<$ & $40$ & $14$ & $100$ & $142$ & $84$\\
\bottomrule
\end{tabular}
\end{table}

\subsection{Proof}

In the proof we multiply by a power of $\rho$ to move $|\sigma(v)|$ into the interval $[1,\Lambda)$, and show that this interval contains no relative unit other than $\pm1$.

\begin{proof}[Proof of Proposition~\ref{prop:units}]
Since $\Lambda>1$, we can multiply any relative unit $v$ by an integer power of $\rho$ so that $1\le a=|\sigma(v)|<\Lambda$.
If $\gamma=0$, then $\alpha^2=1$, and $\alpha\in\Z[\omega]$ gives $v=\pm1$.
Assume $\gamma\ne0$. We derive a contradiction.

Put $v'=\rho\,\tau(v)$. Then $b=|\sigma(v')|=\Lambda/a$, so $b>1$ and $ab=\Lambda$.
The $\theta$-component of $v'$ is also nonzero (if it were zero, then $v=\pm\rho$, contrary to $a<\Lambda$).
Apply Lemma~\ref{lem:gap} to $v$ and $v'$. Combined with $(a-a^{-1})^2\le a^2$, which holds for $a,b\ge1$, this gives $a^2\ge4M$ and $b^2\ge4M$.
Therefore $\Lambda=ab\ge4M$.
For $d=2,3,11,66$, Table~\ref{tab:gap} gives $\Lambda<4M$, a contradiction.

For $d=1$, $\Lambda>32$, so one more step is needed.
From $b^2\ge32$ we get $a^2=\Lambda^2/b^2<50$.
For the $\Z[\omega]$-component $\alpha$ of $v$, we have $\sigma(\alpha)=(\sigma(v)+\sigma(v)^{-1})/2$, so $\sigma(\alpha)^2<14$.
Also $\alpha'^2\le1$.
From $(\sigma(\alpha)+2\alpha')^2\ge0$ we get
\[
(\sigma(\alpha)-\alpha')^2\le\tfrac32\sigma(\alpha)^2+3\alpha'^2<24.
\]
On the other hand, writing $\alpha=h+c\sqrt6$, the left-hand side is $24c^2$. Therefore $c=0$.
Then $\sigma(\alpha)^2=\alpha'^2\le1$. On the other hand, $\sigma(\alpha)^2=1+D_+\sigma(\gamma)^2$ and \eqref{eq:gamma} give $\sigma(\alpha)^2\ge1+D_+D_-=9$, a contradiction.
\end{proof}

This argument does not use the general unit theorem. It needs only the reduction in the real numbers and the rational upper bounds in Table~\ref{tab:gap}.

%======================================================================
\section{Signs and the second initial element}\label{sec:signs}
%======================================================================

By the previous two sections, we can write $\xi=\pm p\rho^m$, where $p$ is an initial element or its relative conjugate and $m\in\Z$.
The map $\tau$ preserves $f$ and the $0$th coordinate, and satisfies $\tau(\rho)=\rho^{-1}$.
Therefore an element of the form $\tau(\xi_i)\rho^m=\tau(\xi_i\rho^{-m})$ may be replaced by $\xi_i\rho^{-m}$.
In what follows we examine $\xi=\pm\xi_i\rho^m$.

\subsection{Nonnegativity}

\begin{lemma}\label{lem:positive}
For each initial element $\xi_i$ and every $m\in\Z$, $P(\xi_i\rho^m)\ge0$.
\end{lemma}

\begin{proof}
All coefficients of the multiplication tables of the five rings are nonnegative. Therefore the set $C$ of elements whose coordinates are all nonnegative is closed under multiplication. The element $\rho$ also belongs to $C$.

For each initial element, we check directly that $\xi_i\rho^2$ and $\tau(\xi_i)\rho^2$ belong to $C$.
Then for $m\ge2$, $\xi_i\rho^m=(\xi_i\rho^2)\rho^{m-2}\in C$.
For $m\le-2$, $\tau(\xi_i\rho^m)=\tau(\xi_i)\rho^{-m}\in C$, and $\tau$ preserves the $0$th coordinate.
The remaining cases $m=-1,0,1$ are checked directly.
\end{proof}

Since $P(\xi)=n+3>0$, the form $\xi=-\xi_i\rho^m$ does not occur.

\subsection{The second initial element for \tsec{$d=2$}{d=2}}

\begin{lemma}\label{lem:extra}
For $d=2$, $\xi=\xi_1\rho^m$ does not occur.
\end{lemma}

\begin{proof}
Multiplication by $\rho^2=(41,24,6,4)$ does not change the first two coordinates modulo $8$. This can be checked for an arbitrary element from the multiplication table.
If $m$ is even, the first two coordinates of $\xi_1\rho^m$ are $(2,3)$ modulo $8$. Then $f\equiv8\pmod{16}$, contrary to $f(\xi)=0$.
If $m$ is odd, the $0$th coordinate is $4$ modulo $8$. But $P(\xi)=2r^2$, and a square is never $2$ modulo $4$. This is a contradiction.
\end{proof}

Thus the only elements left to examine are $\xi_0\rho^m$ ($m\in\Z$).

%======================================================================
\section{The \tsec{$2$}{2}-adic uniqueness lemma}\label{sec:dyadic}
%======================================================================

In this section we prove a general lemma that derives $i=0$ from $f(y\,W^i)=0$.
Here $y$ is a zero of $f$, and $W$ is an invertible element that is $2$-adically close to $1$.
We do not use the zero theorems of $p$-adic analysis. We prove the lemma by repeated divisibility by powers of $2$.

In what follows, a zero $y$ of $f$ to which we apply this lemma is called a \textbf{base point}.
When we write $W=1+2qU$, we call $U$ the \textbf{direction} of $W$.

\subsection{Expansion}

For $y,z\in R$ and $c\in\Z$,
\begin{equation}\label{eq:expand}
f(y+cz)=f(y)+c\lambda_y(z)+c^2H(z),
\end{equation}
where
\[
\lambda_y(z)=2\bigl\{(P(y)-3)P(z)-eQ(y)Q(z)\bigr\},\qquad H(z)=P(z)^2-eQ(z)^2.
\]
The map $\lambda_y$ is $\Z$-linear in $z$ and depends only on $P(z),Q(z)$.

In what follows, write the coordinates of $y$ as $(y_0,y_1,y_2,y_3)$.

\subsection{Expansion of powers}

\begin{lemma}\label{lem:power}
Let $t\ge0$, and suppose that $W=1+2^{t+1}U$ is invertible. Then for every $k\ge0$ and $\ell\in\Z$ the following hold.
\begin{enumerate}
\item There is $E\in R$ such that $W^{2^k\ell}=1+2^{k+t+1}\bigl\{\ell U+2^tU^2E\bigr\}$.
\item If $t\ge1$, there are $b\in\Z$ and $E\in R$ such that $W^{2^k\ell}=1+2^{k+t+1}\bigl\{\ell U+2^tU^2(b+2E)\bigr\}$.
\end{enumerate}
\end{lemma}

\begin{proof}
First, we expand $W^{2^k}$ by repeated squaring, by induction on $k$.
Next, for $\ell\ge0$ we expand the powers of $W^{2^k}$ by the binomial theorem, and for $\ell<0$ we use the expansion of the same form for the inverse $W^{-2^k}$.
In (2), we not only keep the error term as a multiple of $U^2$, but also keep the integer part $b$ of its coefficient separate.
In this way, the form $b+2E$ survives repeated squaring.
\end{proof}

In what follows, put $q=2^t$. The form in (2) can be written as $W^{2^k\ell}=1+2^{k+1}q\{\ell U+qU^2(b+2E)\}$.

\subsection{The uniqueness lemma}

\begin{lemma}[$2$-adic uniqueness lemma]\label{lem:unique}
Let $q=2^t$ ($t\ge1$), let $W=1+2qU$ be an invertible element, and let $y\in R$.
Let $U^2=hG$ with $h\in\{2,4\}$ and $G\in R$, and put $D=2qh^2$. Assume the following.
\begin{enumerate}
\item $f(y)=0$.
\item $h\mid P(yU)$ and $h\mid Q(yU)$.
\item Either $h=2$, or $h=4$ and $\chi(x)=P(x)+Q(x)\bmod2$ is a multiplicative map $R\to\F_2$ with $\chi(G)=0$.
\item $2h^2\mid\lambda_y(yU^2)$.
\item $\lambda_y(yU)\equiv qh^2\pmod D$.
\end{enumerate}
Then, for $i\in\Z$, $f(yW^i)=0$ implies $i=0$.
\end{lemma}

\begin{proof}
First, we show from (1) and (3) that, for every $E\in R$,
\begin{equation}\label{eq:anyerror}
h^2\mid\lambda_y(yU^2E).
\end{equation}
If $h=2$, this is clear, since $\lambda_y$ is a multiple of $2$ to begin with and $U^2=2G$.
If $h=4$, use $f(y)\equiv0\pmod2$. A square is congruent to the number itself modulo $2$, so $P(y)-3\equiv eQ(y)\pmod2$. Therefore
\[
\lambda_y(z)/2=(P(y)-3)P(z)-eQ(y)Q(z)\equiv(P(y)-3)\chi(z)\pmod2.
\]
From $\chi(yGE)=\chi(y)\chi(G)\chi(E)=0$ we get $4\mid\lambda_y(yGE)$. Together with $U^2=4G$, this gives \eqref{eq:anyerror}.

Next, we show for $k\ge0$ that $2^k\mid i$ implies $2^{k+1}\mid i$.
Once this is shown, $i$ is divisible by every power of $2$, so $i=0$.

Write $i=2^k\ell$. By Lemma~\ref{lem:power}~(2), $yW^i=y+cZ$ with $c=2^{k+1}q$ and $Z=y\{\ell U+qU^2(b+2E)\}$.
By \eqref{eq:expand} and $f(y)=0$, we have $0=f(yW^i)=c\{\lambda_y(Z)+cH(Z)\}$, and so $\lambda_y(Z)+cH(Z)=0$.

Here $\lambda_y(Z)=\ell\lambda_y(yU)+qb\,\lambda_y(yU^2)+2q\lambda_y(yU^2E)$.
By \eqref{eq:anyerror} and (4), (5),
\[
\lambda_y(Z)\equiv\ell qh^2\pmod{2qh^2}.
\]
On the other hand, by (2) and $U^2=hG$, $P(Z)$ and $Q(Z)$ are multiples of $h$, and $H(Z)$ is a multiple of $h^2$. Since $c$ is a multiple of $2q$, $cH(Z)\equiv0\pmod{2qh^2}$.
Hence $\ell qh^2\equiv0\pmod{2qh^2}$, that is, $\ell$ is even, and $2^{k+1}\mid i$.
\end{proof}

In what follows, we show that hypotheses (2)--(5) follow in general from the arithmetic of the relative norm of the direction $U$ and from congruences along the orbit.
At individual base points we check only (1).

\subsection{The relative norm of the direction}

\begin{lemma}\label{lem:direction}
If $W=1+4jU$ ($j\in\Z$, $j\ne0$) is a relative unit, then
\[
\pi(U)=-2j\Nm(U),\qquad U^2=-W\Nm(U).
\]
\end{lemma}

\begin{proof}
We have $\Nm(W)=1+4j(U+\tau U)+16j^2\Nm(U)=1$ and $U+\tau U=2\pi(U)$. This gives the first equation.
The second follows from $U^2=U(U+\tau U)-U\tau U=-4jU\Nm(U)-\Nm(U)$.
\end{proof}

An element $W=1+2qU$ with $q\ge2$ has the form of this lemma ($j=q/2$).

In each ring, the contribution of the products of the $\theta$-coefficients to $\pi$ is a multiple of $h$, where $h$ takes the following values:
\[
h=2\quad(d=1,3,11),\qquad h=4\quad(d=2,66).
\]
In fact, for $d\ne11$ both coefficients of $\Delta$ are divisible by $h$. For $d=11$ the basis contains $\omega\theta/11$, so we check this directly with formulas with integer coefficients.

\begin{corollary}\label{cor:width}
For $W$ as in Lemma~\ref{lem:direction}, we have $\Nm(U)\in h\Z[\omega]$, $\pi(U)\in2h\Z[\omega]$ and $U^2\in hR$.
Moreover, for every $y\in R$, $P(yU)$ and $Q(yU)$ are multiples of $h$.
\end{corollary}

\begin{proof}
Since $\pi(U)=-2j\Nm(U)$ has even coefficients, $\pi(U)^2\in4\Z[\omega]$.
The relative norm $\Nm(U)$ is the sum of $\pi(U)^2$ and the contribution of the products of the $\theta$-coefficients, so $\Nm(U)\in h\Z[\omega]$. Then $\pi(U)=-2j\Nm(U)\in2h\Z[\omega]$.
From $U^2=-W\Nm(U)$ we get $U^2\in hR$.
Finally, $\pi(yU)$ is the sum of $\pi(y)\pi(U)$ and the contribution of the products of the $\theta$-coefficients, so it is a multiple of $h$.
\end{proof}

This gives (2) of Lemma~\ref{lem:unique}, and determines $G=U^2/h$.

\subsection{Hypothesis (3)}

In the two cases with $h=4$ ($d=2,66$), $e=3,11$ is odd and $\Delta$ has even coefficients.
Then $\chi(x)=x_0+x_1\bmod2$ is nothing but the substitution $\omega\mapsto1$, $\theta\mapsto0$.
This substitution is compatible with $\omega^2=e\equiv1$ and $\theta^2=\Delta\equiv0\pmod2$, so $\chi$ is a ring homomorphism $R\to\F_2$.

Next, we show $\chi(G)=0$. Since $\chi$ depends only on the $\pi$-part, it suffices to read the $\pi$-part of $G=U^2/4$.
By Corollary~\ref{cor:width}, $\pi(U)^2/4$ has even coefficients.
The contribution of the products of the $\theta$-coefficients is a multiple of $\Delta/4=5+3\omega$ or $189+57\omega$.
Both vanish under $\chi$, so $\chi(G)=0$.

\subsection{Congruences along the orbit}

To check hypotheses (4) and (5) we do not use that the base point $y$ is a zero of $f$.
Instead we use congruences modulo $4$ that hold for every element of the orbit $\xi_0\rho^i$.
Write $\Nm(x)=H(x)+K(x)\omega$ with $H(x),K(x)\in\Z$.

\begin{lemma}\label{lem:orbit4}
For every $i\in\Z$, the element $y=\xi_0\rho^i$ satisfies $f(y)\equiv0$ and $K(y)\equiv0\pmod4$.
\end{lemma}

\begin{proof}
Reading the multiplication table of each ring modulo $4$, the products of the $\theta$-coefficients contribute only multiples of $4$ to $H$.
Hence $H(y)\equiv y_0^2+ey_1^2\pmod4$, and
\[
f(y)+H(y)-11\equiv2y_0(y_0-3)\equiv0\pmod4,
\]
because $y_0$ and $y_0-3$ have different parities.
Every element of the orbit satisfies $\Nm(y)=\Nm(\xi_0)\Nm(\rho)^i=11$, so $H(y)=11$ and $K(y)=0$, and the claim follows.
\end{proof}

\subsection{Hypothesis (4)}

Put $\Nm(U)=u+v\omega$. Corollary~\ref{cor:width} gives $h\mid u$ and $h\mid v$; in fact $4\mid u$ in every ring,
because the products of the $\theta$-coefficients contribute a multiple of $4$ to the constant term.

By Lemma~\ref{lem:direction}, $yU^2=-y\Nm(U)-4j(yU)\Nm(U)$.
In the second term the $\pi$-part is a multiple of $h^2$ by Corollary~\ref{cor:width}, so its value under $\lambda_y$ is a multiple of $8jh^2$.
For the first term, for every $y$,
\begin{equation}\label{eq:scalar}
\lambda_y(y\Nm(U))=2\bigl\{[y_0(y_0-3)-ey_1^2]u-3ey_1v\bigr\}.
\end{equation}
So it suffices to show that the expression in braces is a multiple of $h^2$.

If $h=2$, then $4\mid u$, $2\mid v$ and $2\mid e$, so the braces are a multiple of $4$ whatever $y$ is.

If $h=4$, use $y_0(y_0-3)-ey_1^2=f(y)+3y_0-11$ to rewrite \eqref{eq:scalar} as
\[
\lambda_y(y\Nm(U))=2\bigl\{(3y_0-11)u-3ey_1v\bigr\}+2u\,f(y).
\]
By Lemma~\ref{lem:orbit4}, $4\mid f(y)$, so the last term is a multiple of $32$.
For the rest we use two congruences. Put $\varsigma=-1$ ($d=2$) and $\varsigma=1$ ($d=66$).
\begin{itemize}
\item Orbit congruence: for every $m\in\Z$, $y=\xi_0\rho^m$ satisfies $y_0+\varsigma y_1\equiv1\pmod4$.
In fact, for every $x$ the difference $(P+\varsigma Q)(x\rho)-(P+\varsigma Q)(x)$ is a multiple of $4$, and the value at $\xi_0$ is $1$.
\item Direction congruence: $u-\varsigma v\equiv0\pmod{16}$.
This follows from $\pi(U)\in8\Z[\omega]$ and the expansion of $\Nm(U)$ with the $\theta$-coefficient written as $c+k\omega$.
\end{itemize}
Since $e\equiv-1\pmod4$, we have $(3y_0-11)-3e\varsigma y_1=3(y_0+\varsigma y_1-1)-8-3\varsigma(e+1)y_1\equiv0\pmod4$.
Together with $v\equiv\varsigma u\pmod{16}$ and $4\mid u$, the braces are a multiple of $16$.
Hence $32\mid\lambda_y(y\Nm(U))$ when $h=4$, and hypothesis (4) holds.

\subsection{Hypothesis (5): the case \tsec{$h=2$}{h=2}}

In the three cases with $h=2$ ($d=1,3,11$), the multiplication tables give
\begin{equation}\label{eq:Kmod4}
K(y)\equiv2(y_0y_1-y_2^2)\pmod4
\end{equation}
for every $y$.
Indeed, $K(y)/2$ expands to $y_0y_1-9y_2^2-44y_2y_3-54y_3^2$, $y_0y_1-3y_2^2-6y_3^2$ and $y_0y_1-11y_2^2-16y_2y_3-6y_3^2$, and modulo $2$ only $y_0y_1-y_2^2$ remains.

The congruences $f(y)\equiv K(y)\equiv0\pmod4$ of Lemma~\ref{lem:orbit4} show that $y_0,y_1,y_2$ are all odd.
Indeed, since $e\equiv2\pmod4$, $f(y)\equiv0$ reads $(y_0-3)^2+2y_1^2\equiv2\pmod4$.
Squares are $0$ or $1$ modulo $4$, so $y_1$ is odd and $y_0-3$ is even.
Then \eqref{eq:Kmod4} and $K(y)\equiv0$ give $y_2^2\equiv y_0y_1\equiv1\pmod2$.

By Corollary~\ref{cor:width}, $P(U)$ and $Q(U)$ are multiples of $4$.
Reading the multiplication tables modulo $4$ gives $P(yU)\equiv0$ and $Q(yU)\equiv2y_2U_2\pmod4$, where $U_2$ is the coordinate of index $2$ of $U$.
Hence, if $U_2$ is odd,
\[
\lambda_y(yU)=2(y_0-3)P(yU)-2ey_1Q(yU)\equiv-8y_1y_2U_2\equiv8\pmod{16},
\]
which is hypothesis (5) for $q=2$, $h=2$.
For the three directions used, $U_2=9,99,57$, all odd.

\subsection{Hypothesis (5): the case \tsec{$h=4$}{h=4}}

In the two cases with $h=4$, write $\Delta=4(k+\ell\omega)$; then $(e,k,\ell)$ is $(3,5,3)$ for $d=2$ and $(11,189,57)$ for $d=66$.
Let $W=1+4jU$, and write the product of the $\theta$-coefficients $U_2+U_3\omega$ of $U$ and $y_2+y_3\omega$ of $y$ as $X+Y\omega$, that is,
\[
X=y_2U_2+ey_3U_3,\qquad Y=y_2U_3+y_3U_2.
\]
By the argument for hypothesis (4), the part $\pi(U)=-2j\Nm(U)$ contributes only multiples of $64j$ to $\lambda_y$.
Computing the contribution of the $\theta$-part gives
\begin{equation}\label{eq:pairing}
\lambda_y(yU)\equiv8\bigl([k(y_0-3)-e\ell y_1]X+e[\ell(y_0-3)-ky_1]Y\bigr)\pmod{64j}.
\end{equation}

The coefficients on the right depend only on the residues of the coordinates of $y$ modulo small numbers.
We therefore find the residue classes containing the base points as congruences preserved along the orbit.
Write $\rho^4=1+4S$. By Corollary~\ref{cor:width}, the first two coordinates of $S$ are multiples of $8$, and for every $z$ the first two coordinates of $zS$ are multiples of $4$.
Hence $z\rho^4-z=4zS$ has its first two coordinates divisible by $16$ and the other two by $4$.
In other words, the residues of the coordinates modulo $(16,16,4,4)$ do not change under multiplication by $\rho^{\pm4}$.
For $d=2$ and the class $m\equiv0\pmod4$, $S_3$ and $z_0$ are even, so the coordinate of index $2$ is preserved modulo $8$ as well.
The resulting classes are given in Table~\ref{tab:phases}. The starting points are $\xi_0$, $\xi_0\rho^{-1}=(8,3,-2,0)$ ($d=2$) and $\xi_0\rho=(66,19,0,1)$ ($d=66$).

\begin{table}[htbp]
\centering
\caption{The residue classes containing the base points for $h=4$ ($m$ modulo $4$), and the value of \eqref{eq:pairing}}\label{tab:phases}
\footnotesize\setlength{\tabcolsep}{2.5pt}
\begin{tabular}{cccccc}
\toprule
case & moduli & residue of $y$ & $W$ & $\lambda_y(yU)$ & condition\\
\midrule
$d=2$, $m\equiv0$ & $(16,16,8,4)$ & $(2,1,2,3)$ & $\rho^8$ & $16(6U_2+U_3)\bmod128$ & $6U_2+U_3\equiv4\ (8)$\\
$d=2$, $m\equiv-1$ & $(16,16,4,4)$ & $(8,3,2,0)$ & $\rho^4$ & $32U_2\bmod64$ & $U_2$ odd\\
$d=66$, $m\equiv1$ & $(16,16,4,4)$ & $(2,3,0,1)$ & $\rho^4$ & $-16U_3\bmod64$ & $U_3\equiv2\ (4)$\\
\bottomrule
\end{tabular}
\end{table}

When the condition in the last column holds, $\lambda_y(yU)$ is $64\bmod128$, $32\bmod64$ and $32\bmod64$, respectively, which is hypothesis (5).
For the two directions of $d=2$ it suffices to check, for the short direction $S$ ($\rho^4=1+4S$), that $S_2$ is odd and $6S_2+S_3\equiv4\pmod8$.
Indeed, the long direction is $U=(\rho^8-1)/8=S+2S^2$, and $S^2\in4R$ gives $U\equiv S\pmod8$.

Thus hypotheses (2)--(5) all hold in general.
The finite facts checked individually are only that $f(y)=0$ at the nine base points, the six conditions on the directions, and the coordinates of the two starting points.

%======================================================================
\section{Determination of the exponent}\label{sec:exponents}
%======================================================================

The base points $y$ and the units $W$ at which we use Lemma~\ref{lem:unique} are given in Table~\ref{tab:centres}.
All base points have the form $\xi_0\rho^{m_0}$, and $f(y)=0$ is checked directly.

\begin{table}[htbp]
\centering
\caption{The base points at which the uniqueness lemma is used}\label{tab:centres}
\begin{tabular}{cccccc}
\toprule
$d$ & exponent $m_0$ of the base point & $W$ & $q$ & $h$ & $D$\\
\midrule
$1$ & $0,\ -1$ & $\rho^2$ & $2$ & $2$ & $16$\\
$2$ & $0,\ -4$ & $\rho^8$ & $4$ & $4$ & $128$\\
$2$ & $-1$ & $\rho^4$ & $2$ & $4$ & $64$\\
$3$ & $0,\ -1$ & $\rho^2$ & $2$ & $2$ & $16$\\
$11$ & $0$ & $\rho^2$ & $2$ & $2$ & $16$\\
$66$ & $1$ & $\rho^4$ & $2$ & $4$ & $64$\\
\bottomrule
\end{tabular}
\end{table}

That $W$ has the form $1+2qU$ follows from the behaviour of $\rho$ modulo $4$.
For $d=1,3,11$, $\rho\equiv1\pmod2$, so $\rho^2\equiv1\pmod4$.
For $d=2,66$, $\rho^2\equiv1+2\theta\pmod4$, and its square $\rho^4$ is $\equiv1\pmod4$.
Moreover, if $\rho^4=1+4U'$, then $\rho^8=1+8(U'+2U'^2)$, so $\rho^8$ for $d=2$ has the form with $q=4$.

In what follows, we determine the $m$ with $f(\xi_0\rho^m)=0$.

\paragraph{$d=1,3$.}
Write $m=2a$ if $m$ is even and $m=-1+2a$ if $m$ is odd, and use Lemma~\ref{lem:unique} at the base points $\xi_0$ and $\xi_0\rho^{-1}$, respectively.
Therefore $m=0$ or $m=-1$.

\paragraph{$d=11$.}
Since $\tau(\xi_0)=\xi_0\rho^{-1}$, we have $\tau(\xi_0\rho^m)=\xi_0\rho^{-1-m}$.
Since $\tau$ preserves $f$ and the $0$th coordinate, we may replace $m$ by $-1-m$ and assume that $m$ is even.
Using Lemma~\ref{lem:unique} at the base point $\xi_0$, we get $m=0$.

\paragraph{$d=66$.}
First, we show from the original square condition that $m$ is odd. This step does not use $f$.
Since $\xi_0=\omega$, if $m=2i$, then $P(\xi)=11\,Q(\rho^{2i})=66r^2$, that is, $Q(\rho^{2i})=6r^2$.
Write $\rho^2=1+2S$ with $S=(756,228,19,6)$, and let $i=2^ka$.
Applying Lemma~\ref{lem:power}~(1) to $W=\rho^2$, $t=0$, we get $\rho^{2i}=1+cV$ with $c=2^{k+1}$ and $V=aS+S^2E$.
Compare the $\omega$-coefficients of $\Nm(\rho^{2i})=1$, and use that the $\omega$-coefficient of the relative norm of $V$ is congruent to $12a^2$ modulo $48$. Then we get an integer $g$ with
\[
r^2=-c^2(a^2+4g).
\]
Since $c^2\mid r^2$, we can write $r=cs$. Then $s^2+a^2\equiv0\pmod4$, so $a$ is even.
Repeating this argument for $k=0,1,2,\dots$, $i$ is divisible by every power of $2$. Hence $i=0$, and $Q(\rho^0)=0=6r^2$ contradicts $r>0$.
Therefore $m$ is odd.

Next, since $\tau(\xi_0)=\xi_0$, we have $\tau(\xi_0\rho^m)=\xi_0\rho^{-m}$.
Replacing $m$ by $-m$, we may assume $m\equiv1\pmod4$.
Using Lemma~\ref{lem:unique} at the base point $\xi_0\rho$, we get $m=1$.

\paragraph{$d=2$.}
This case has two close roots $0$ and $-4$, so one preparatory step is needed.
Every element $x=\xi_0\rho^m$ of the orbit satisfies $2\mid x_0$ and $x_1+2x_3\equiv3\pmod4$.
For every $x$ satisfying this condition, expanding the multiplication gives
\[
f(x\rho^3)-3f(x\rho^2)+3f(x\rho)-f(x)\equiv0\pmod{128}.
\]
From the first three values and induction in both directions using this relation, we get, for every $m\in\Z$,
\begin{equation}\label{eq:entrance}
f(\xi_0\rho^m)\equiv32m(m+1)\pmod{128}.
\end{equation}
Therefore, if $f(\xi_0\rho^m)=0$, then $m\equiv0$ or $m\equiv-1\pmod4$.
If $m=4a$, we split by the parity of $a$ and use Lemma~\ref{lem:unique} at the base points $\xi_0$ and $\xi_0\rho^{-4}$ with $W=\rho^8$. We get $m=0$ or $m=-4$.
If $m=-1+4a$, we use the base point $\xi_0\rho^{-1}$ with $W=\rho^4$ and get $m=-1$.

%======================================================================
\section{Proof of the main theorem}\label{sec:conclusion}
%======================================================================

\begin{proof}[Proof of Theorem~\ref{thm:main}]
By the above, $\xi$ is one of the nine elements in Table~\ref{tab:values}.
Since $P(\xi)=n+3$, the $0$th coordinate determines $n$.

\begin{table}[htbp]
\centering
\caption{The remaining elements and $n$}\label{tab:values}
\begin{tabular}{cccc}
\toprule
$d$ & $m$ & $P(\xi_0\rho^m)$ & $n$\\
\midrule
$1$ & $0$ & $1$ & $-2$\\
$1$ & $-1$ & $25$ & $22$\\
$2$ & $0$ & $2$ & $-1$\\
$2$ & $-4$ & $13778$ & $13775$\\
$2$ & $-1$ & $8$ & $5$\\
$3$ & $0$ & $3$ & $0$\\
$3$ & $-1$ & $27$ & $24$\\
$11$ & $0$ & $11$ & $8$\\
$66$ & $1$ & $66$ & $63$\\
\bottomrule
\end{tabular}
\end{table}

The negative values $n=-2,-1$ contradict $n\in\N$. The remaining seven are the values in the theorem. The largest solution $13775$ comes from $m=-4$ in the case $d=2$.
Conversely, substitution shows that the seven values are solutions with $m=1,6,11,45,51,209,660099$, respectively.
\end{proof}


%======================================================================
\appendix
\section{Formalization in Lean}\label{app:lean}
%======================================================================

The final theorem has the following form. The division is carried out in the rationals $\Q$.
{\small
\begin{verbatim}
theorem Tetrahedral.tetrahedral_add_one_eq_sq_iff (n : ℕ) :
    (∃ m : ℕ, (n * (n + 1) * (n + 2) / 6 + 1 : ℚ) = m ^ 2) ↔
      n = 0 ∨ n = 5 ∨ n = 8 ∨ n = 22 ∨ n = 24 ∨ n = 63 ∨ n = 13775
\end{verbatim}
}

The correspondence between the sections of the text and the Lean files is given in Table~\ref{tab:lean}.
The paths in the table are relative to the directory \texttt{lean/Tetrahedral/}. The extension \texttt{.lean} is omitted.

\begin{table}[htbp]
\centering
\small
\caption{Correspondence between the text and Lean}\label{tab:lean}
\begin{tabular}{lL{10.4cm}}
\toprule
Section & Files\\
\midrule
\ref{sec:branches} & \texttt{Basic}, \texttt{Squarefree}, \texttt{IntegerParity}, \texttt{IntegerSquareObstructions}, \texttt{Branches}\\
\ref{sec:parameters} & \texttt{SpecialPell}, \texttt{QuadraticParameters}, \texttt{SelectedParameters}, \texttt{SquareFactor}, \texttt{D3Parameters}, \texttt{NormEleven/Parameters}\\
\ref{sec:norm11} & \texttt{Quartic/D*Ring}, \texttt{Quartic/D1}, \texttt{D2}, \texttt{D3}, \texttt{D11}, \texttt{D66}, \texttt{NormEleven/Certificates}\\
\ref{sec:division} & \texttt{Quartic/D1Division}, \texttt{Quartic/D2Division}, \texttt{Quartic/D3Division}, \texttt{Quartic/D11}, \texttt{Quartic/D66}\\
\ref{sec:units} & \texttt{Quartic/UnitGap}, \texttt{Quartic/RelativeUnitGap}, \texttt{Quartic/D*Trace}\\
\ref{sec:signs} & \texttt{PositiveOrbit}, \texttt{OrbitCongruence}, \texttt{NormEleven/PositiveOrbitData}, \texttt{NormEleven/D2Congruences}\\
\ref{sec:dyadic} & \texttt{Binomial}, \texttt{QuadraticPolynomial}, \texttt{FactoredExponent}, \texttt{NormDerivative}, \texttt{RelativeDirection}, \texttt{NormEleven/RelativeDirectionData}, \texttt{NormEleven/ParityDirectionData}\\
\ref{sec:exponents} & \texttt{OrbitReflection}, \texttt{FiniteDifference}, \texttt{NormEleven/DyadicSimple}, \texttt{NormEleven/DyadicCertificates}, \texttt{NormEleven/D2Entrance}, \texttt{NormEleven/D2Dyadic}, \texttt{NormEleven/D11Dyadic}, \texttt{NormEleven/D66Square}, \texttt{NormEleven/D66Dyadic}\\
\ref{sec:conclusion} & \texttt{NormEleven/D1}, \texttt{D2}, \texttt{D3}, \texttt{D11}, \texttt{D66}, \texttt{Classification}\\
\bottomrule
\end{tabular}
\end{table}

Some symbols and names differ between the text and Lean. The main ones are listed in Table~\ref{tab:names}.

\begin{table}[htbp]
\centering
\small
\caption{Symbols in the text and in Lean}\label{tab:names}
\begin{tabular}{ll}
\toprule
Text & Lean\\
\midrule
$\omega=\sqrt e$ & \texttt{q}\\
$\theta$ & \texttt{s} (\texttt{t} for $d=66$)\\
the two coefficients of $\Nm(x)$ & \texttt{relH x}, \texttt{relK x}\\
$\tau$ & \texttt{tau}\\
$\sigma$ & \texttt{sig}\\
base point & centre\\
direction $U$ & \texttt{step}\\
$h$ & width\\
$\chi$ & parity character\\
\bottomrule
\end{tabular}
\end{table}

The following three files are proofs generated by the Python scripts in \texttt{scripts/}.
\begin{itemize}
\item \texttt{NormEleven/PositiveOrbitData}
\item \texttt{NormEleven/RelativeDirectionData}
\item \texttt{NormEleven/ParityDirectionData}
\end{itemize}
The generated files are also ordinary Lean proofs, and the kernel checks them.
The scripts are kept so that one can check that the same files are regenerated.

\begin{thebibliography}{9}
\bibitem{GPZ94} J.~Gebel, A.~Peth\H{o}, H.~G.~Zimmer, Computing integral points on elliptic curves, \emph{Acta Arith.} \textbf{68} (1994), 171--192.
\bibitem{ST94} R.~J.~Stroeker, N.~Tzanakis, Solving elliptic diophantine equations by estimating linear forms in elliptic logarithms, \emph{Acta Arith.} \textbf{67} (1994), 177--196.
\bibitem{OEIS} The On-Line Encyclopedia of Integer Sequences, A000292 and A216269, \url{https://oeis.org}.
\end{thebibliography}

\end{document}
